\documentclass[11pt,a4paper]{article}
\usepackage{gtu-el-yazisi}
\begin{document}
\notbilgi{Görüntü İşleme}{Morphological Image Processing}{Ders-10.pdf}
\textit{Not: Orijinal dosyadaki basılı sayfalar (s. 3) dijitalleştirilmedi; yalnızca el yazısı sayfalar aktarıldı.}

\sayfa{1}
$\star$ eccentiricity

Alan bazlı filtreleme

\section{Morphological Image Processing}
$\hookrightarrow$ binary

$\hookrightarrow$ grayscale

\[ \mathbb{Z}^2 \quad\leftarrow\quad \text{iki boyutlu tam sayı uzayı} \]

\[ \mathbb{Z}^2_{x,y} \qquad \text{($x,y$ noktası)} \]

\begin{align*}
A &= \{(a_1,b_1),(a_2,b_2),\dots\}\\
  &= \{(x_1,y_1),(x_2,y_2),\dots\}
\end{align*}
\[ a = (a_1,a_2) \]

\[ a \in A \]

\[ a \notin A \qquad \text{(elemanı değil)} \]

\[ A \subseteq B \qquad \text{(A, B'nin alt kümesi)} \]

\begin{gather*}
C = A \cup B \qquad \text{(birleşimi)}\\
D = A \cap B \qquad \text{(ortak elemanlar)}\\
A \cap B \neq \emptyset \qquad \text{(A ve B ayrık)}\footnotemark
\end{gather*}
\footnotetext{Claude notu: ayrık kümeler için muhtemelen $A \cap B = \emptyset$ olmalı.}

\[ A^c = \{w \mid w \notin A\} \qquad \text{(w öyleki A'da değil)} \]

\[ A - B = \{w \mid w \in A,\ w \in B\footnotemark\} = A \cap B^c \]
\footnotetext{Claude notu: muhtemelen $w \notin B$ olmalı.}

\[ \hat{B} = \{w \mid w = -b,\ \text{for } b \in B\} \qquad \text{(reflection)} \]

\[ (A)_z = \{w \mid w = a + z,\ \text{for } a \in A\} \]
A kümesini z noktasına translation

\sayfa{2}
\section{Dilation operator}
$A$ and $B$ sets in $\mathbb{Z}^2$
\[ C = A \oplus \underset{\text{structuring element}}{B} = \{z \mid (\hat{B})_z \cap A \neq \emptyset\} \]
A dilayted by B

\elcizimi[0.7]{sekiller/Ders-10-s2-1.png}[2]

A'yı d kadar büyüttüm

dilate ettim

\section{Erosion:}
\[ C = A \ominus B = \{z \mid (B)_z \subseteq A\} \]

\elcizimi[0.45]{sekiller/Ders-10-s2-2.png}[2]

\[ (A \ominus B)^c = A^c \oplus \hat{B} \]
\[ A \ominus B = (A^c \oplus \hat{B})^c \]

\sekilnotu[2]{Bu sayfa basılı materyal içeriyor (ör. ders slaytı); dijitalleştirilmedi.}

\sayfa{4}
\section{Opening and Closing}
\[ \underline{\text{opening}} \qquad A \circ B = (A \ominus B) \oplus B \]
$\swarrow$ küçültüp \qquad $\swarrow$ büyütüyor
\[ \underline{\text{closing}} \qquad A \bullet B = (A \oplus B) \ominus B \]
$\swarrow$ aslında bu görüntünün negatifi
\[ (A \bullet B)^c = (A^c \circ B^c) \]

\elcizimi[0.5]{sekiller/Ders-10-s4-1.png}[4]

opening yaparsak $3\times 3$'ü götürüz

\sekilnotu[4]{Bu sayfa basılı materyal içeriyor (ör. ders slaytı); dijitalleştirilmedi.}
Closing delikleri kapatıyor

$\leftarrow$ opening

\sayfa{5}
\[ A - (A \ominus B) \quad \rightarrow \text{inner edge} \]
\[ (A \oplus B) - A \quad \rightarrow \text{outer edge} \]

\elcizimi[0.5]{sekiller/Ders-10-s5-1.png}[5]

\[ (A \oplus B) - (A \ominus B) \]
$\hookrightarrow$ inner, outer edge

\elcizimi[0.3]{sekiller/Ders-10-s5-2.png}[5]

Hit-OR-Miss Transform

Hit-And-Miss aslında
\[ A \circledast B = \underbrace{(A \ominus X)}_{\text{hit}} \cap \underbrace{[A^c \ominus (W-X)]}_{\text{Miss}} \]
\[ B = (X, W) \]
\[ A \circledast B = (A \ominus B_1) \cap [A^c \ominus B_2] \]
\[ B = (B_1, B_2) \]

\elcizimi[0.85]{sekiller/Ders-10-s5-3.png}[5]

\sayfa{6}
\elcizimi[0.5]{sekiller/Ders-10-s6-1.png}[6]

\elcizimi[0.35]{sekiller/Ders-10-s6-2.png}[6]

\elcizimi[0.55]{sekiller/Ders-10-s6-3.png}[6]

\section{Hole filling}
\[ X_k = (X_{k-1} \oplus B) \cap I^c \qquad k = 1,2,3,\dots \]
\[ \text{if } X_k = X_{k-1} \Rightarrow \text{stop} \]

delikten bir nokta ile başlamak gerekiyo.
\[ X_1 = (X_0 \oplus B) \cap A^c \]

\sekilnotu[6]{Bu sayfa basılı materyal içeriyor (ör. ders slaytı); dijitalleştirilmedi.}

\sayfa{7}
\[ A \oplus (B \oplus C) = (A \oplus B) \oplus C \]
\[ A \ominus \underbrace{(B \oplus C)}_{D} = (A \ominus B) \ominus C \]
\[ A \oplus B = B \oplus A \]
\[ ((A \circ B) \circ B) \circ B = A \circ B \]
\[ (A \circ B) \circ B = A \bullet B \]

\end{document}
